\documentclass[10pt,letterpaper]{article}
\usepackage[hmargin=1in,vmargin=0.75in]{geometry}
\date{}
\usepackage{xurl}
\setlength{\emergencystretch}{2em}
\title{Scoring Detector Faults in an Archived Freeway Replay}
\author{Arjun Vijay Prakash \\ \small ITS World Congress Birmingham 2027 | Preliminary poster
draft}

\begin{document}

\maketitle
\begin{abstract}
Graph reconstruction reduced speed error in a replay of 58 detector outages and delays
on archived southbound I-880 and westbound I-10 readings. It had 4.65 km/h mean absolute
error on 468 scored target instances, compared with 5.86 km/h for latest-reading
persistence. Both methods missed all ten low-speed queue-proxy positives. The replay
records both the estimated speed and which readings had arrived when it was computed.
\end{abstract}

\section*{Arrival-time replay}
Each archived five-minute Caltrans PeMS file has a capture time and station metadata version.
At each decision time, the replay opens only readings that had arrived. Its ledger records
the last valid observation, its age, the neighboring speeds, and the evidence for a detector
check or queue check. A separate status marks stale feeds and missing or unusable readings.

\section*{Injected fault replay}
We applied 30- and 60-minute blackouts and separate 15-minute delivery delays to archived
detector readings. An original reading qualified as a target only when the detector and its
two neighbors had usable reference speeds. The graph fit used a persistence prior built
from the target corridor's arrived readings. Mean absolute error was 4.67 versus 6.05 km/h
for blackouts, and 4.63 versus 5.67 km/h for delayed readings. Fresh readings returned
inside the observation window in 28 cases. The remaining 30 recovery times are censored.

\section*{Operational implications}
The replay supports a demonstration of how to investigate a missing reading. Its fault
cases use the same brief archive, and targets can recur under different fault settings.
Queue accuracy remains unestablished. The reference consists of adjacent low detector
speeds and supplies no operator-confirmed incident labels.

\section*{References}
[1] Caltrans. Performance Measurement System. \url{https://pems.dot.ca.gov/}

[2] Development evaluation and receipts.
\url{https://github.com/ArjunCodess/CREST/blob/main/docs/crest-evaluation.md}
\end{document}
